\documentclass[12pt,a4paper]{article}
\usepackage[utf8]{inputenc}
\usepackage[T2A]{fontenc}
\usepackage[russian]{babel}
\usepackage{amsmath,amssymb,geometry}
\geometry{margin=2cm}
\begin{document}
\section*{Рабочая тетрадь: квадратные уравнения}
\subsection*{Теория. Дискриминант}
Для квадратного уравнения \[
ax^2 + bx + c = 0, \qquad a \neq 0
\]
дискриминант вычисляется по формуле
\[
D = b^2 - 4ac.
\]
\vspace{8mm}
\subsection*{Задача 1. Найдите корни уравнения}
\[
x^2 - 5x + 6 = 0.
\]
\vspace{8mm}
\subsection*{Задача 2. Найдите корни уравнения}
\[
x^2 - 9 = 0.
\]
\vspace{8mm}
\subsection*{Теория. Число корней}
Если $D > 0$, уравнение имеет два различных корня. Если $D = 0$, уравнение имеет один корень. Если $D < 0$, действительных корней нет.
\vspace{8mm}
\subsection*{Задача 3. Определите число действительных корней}
\[
x^2 + 4x + 5 = 0.
\]
\vspace{8mm}
\end{document}
